% !Mode:: "TeX:UTF-8"
% !TEX program  = xelatex
\section{系统模型与问题定义}
\subsection{系统模型}
本文考虑一个多租户的大语言模型推理系统，其中一个共享的基础模型被加载至 GPU 上，而多个 LoRA 适配器用于支持不同任务或用户请求。设适配器集合为 $\mathcal{A}=\{a_1,a_2,\dots,a_N\}$，系统在运行过程中需要处理一系列请求序列 $\mathcal{R}=\{r_1,r_2,\dots\}$，其中每个请求 $r_i$ 关联唯一的适配器 $a(r_i)\in\mathcal{A}$。

由于 GPU 显存资源有限，系统无法同时容纳所有适配器。设 GPU 可用于适配器的显存预算为 $M$，每个适配器 $a$ 占用显存大小为 $\mathrm{mem}(a)$，则任意时刻驻留在 GPU 上的适配器集合 $\mathcal{H}_t\subseteq\mathcal{A}$ 必须满足如下约束：
\begin{equation}
\sum_{a\in\mathcal{H}_t}\mathrm{mem}(a)\le M
\end{equation}

当一个请求到达时，若其所需适配器已经存在于当前驻留集合 $\mathcal{H}_t$ 中，则可以直接执行推理；否则，系统需要将该适配器从 CPU 或外部存储加载至 GPU。该过程通常伴随着显著的时间开销，并可能触发已有适配器的淘汰操作。此外，不同适配器在显存占用与加载开销上可能存在显著差异，这进一步增加了资源管理的复杂性。因此，系统需要在运行过程中持续做出决策，以动态维护适配器驻留集合 $\mathcal{H}_t$。

\subsection{问题定义}
在上述系统模型下，本文关注如下核心问题：在给定请求序列 $\mathcal{R}$ 和显存约束 $M$ 的条件下，如何动态管理适配器的驻留与淘汰，以优化系统整体性能。具体而言，在每个时间步 $t$，系统接收到一个请求 $r_t$，其对应适配器为 $a_t=a(r_t)$。系统需要根据当前状态（包括历史请求、当前驻留集合以及系统负载情况）决定是否需要加载新的适配器，并在必要时选择合适的适配器进行淘汰。该问题可以形式化为一个在线决策问题：系统在未知未来请求的情况下，基于当前及历史信息，对适配器驻留集合 $\mathcal{H}_t$ 进行动态更新，从而最小化长期运行成本。

本文将系统优化目标定义为最小化适配器管理带来的总体开销。设在时间步 $t$ 是否发生适配器加载由指示函数 $\mathbb{I}_{\text{miss}}(t)$ 表示，则总体开销可表示为：
\begin{equation}
\min\sum_t\mathbb{I}_{\text{miss}}(t)\cdot\mathrm{cost}(a_t)
\end{equation}
其中 $\mathrm{cost}(a_t)$ 表示加载适配器 $a_t$ 的开销，其可能与适配器大小、带宽以及系统状态相关。

该优化目标不仅影响平均延迟，还会对尾延迟（如 P95 latency）和系统吞吐量产生重要影响。此外，由于适配器大小不一致，简单地最小化缓存未命中次数并不能直接反映系统性能，因此需要在决策过程中同时考虑适配器的显存占用与加载成本。

\subsection{基线方法}
\subsubsection{FIFO (First-In, First-Out)}
FIFO 是最简单的适配器管理策略之一。在该方法中，适配器按照进入 GPU 的时间顺序进行管理。当系统需要加载新的适配器且显存不足时，将优先淘汰最早加载到 GPU 中的适配器。该方法不利用任何访问历史信息，也不考虑适配器的使用频率或未来需求，因此属于完全被动的策略。

\subsubsection{LRU (Least Recently Used)}
LRU 基于时间局部性原理进行适配器管理。该方法假设近期被访问的适配器在未来更可能被再次使用。因此，当需要淘汰适配器时，系统会优先移除最长时间未被访问的适配器。具体而言，系统为每个适配器维护最近一次访问时间戳。当发生显存不足时，选择时间戳最早的适配器进行淘汰。

\subsubsection{LFU (Least Frequently Used)}
LFU 基于访问频率进行适配器管理。该方法通过统计每个适配器的访问次数，优先淘汰访问频率最低的适配器。LFU 能够较好地识别长期热点适配器，但其主要局限为缺乏时间敏感性，在访问模式发生变化时，LFU 难以及时适应新的热点分布。

\subsubsection{方法对比分析}
上述三种方法均属于被动缓存策略，即仅基于历史访问信息进行决策，而不显式建模未来请求。这些方法在传统缓存场景中具有一定效果，但在多 LoRA 适配器推理环境下存在明显局限，比如未考虑适配器显存占用差异以及不具备对未来需求的预测能力。
